\chapter{Reading the maya source}
\label{ch:readingsource}

\begin{aside}
The library is small enough to read in an afternoon. Here's
the tour.
\end{aside}

\section{Layout}

\begin{itemize}[leftmargin=*]
  \item \file{maya/include/maya/core/}: foundational types.
  \item \file{maya/include/maya/dsl/}: builders and pipe
    operators.
  \item \file{maya/include/maya/element/}: element types.
  \item \file{maya/include/maya/layout/}: flexbox engine.
  \item \file{maya/include/maya/render/}: canvas, diff,
    emit.
  \item \file{maya/include/maya/runtime/}: signals,
    effects, loop.
  \item \file{maya/include/maya/widget/}: catalogue.
  \item \file{maya/include/maya/term/}: ANSI, input parser.
\end{itemize}

\section{Reading order}

\begin{enumerate}[leftmargin=*]
  \item Start with \file{maya/include/maya/maya.hpp}: the
    umbrella header.
  \item Skim \file{element/element.hpp}: the base type.
  \item Read \file{render/canvas.hpp}: packed cells.
  \item Read \file{render/diff.hpp}: the loop.
  \item Read \file{runtime/signal.hpp}: reactivity.
\end{enumerate}

Each is a few hundred lines.

\section{Naming conventions}

\begin{itemize}[leftmargin=*]
  \item snake\_case for functions and variables.
  \item CamelCase for types.
  \item All lowercase for the namespace (\code{maya}).
  \item Trailing underscore on private members.
\end{itemize}

\section{Templates vs concepts}

\maya{} uses \cpp{}20 concepts liberally. A function
constraint says ``this argument must satisfy X''.

\section{Tests}

\file{maya/tests/} mirrors the include layout. Most files
have a corresponding test.

\section{Examples}

\file{maya/examples/} has small programs demonstrating
features. Useful as starting points.

\section{Recap}

\begin{recap}
\begin{itemize}[leftmargin=*]
  \item Source under \file{maya/include/maya/}.
  \item Read element \textrightarrow{} canvas
    \textrightarrow{} diff \textrightarrow{} signal.
  \item Tests mirror the layout.
\end{itemize}
\end{recap}

\section*{Workshop}

\begin{enumerate}[leftmargin=*]
  \item Read \file{element/element.hpp} top to bottom.
  \item Find one widget in \file{widget/} and trace its
    paint method.
  \item Run the examples; modify one.
\end{enumerate}
